\textbf{Время:} 01:30

\textbf{Тезис:} измеренная пионная корреляция содержит и обменное усиление, и кулоновское отталкивание.

\textbf{На экране:} заголовок и отметка «STAR · эксперимент»; затем три панели out, side, long и их оси; после этого красные и синие точки, линии через них и нижняя подпись об учёте кулоновского отталкивания.

\textbf{Текст:} Перед нами уже измеренные данные STAR. Три панели соответствуют только что определённым направлениям out, side и long. По горизонтали — соответствующая компонента разности импульсов, по вертикали — корреляционная функция.

Красные точки относятся к парам положительных пионов, синие — к парам отрицательных. Показан наиболее центральный класс, 0–10 процентов. Здесь выбраны пары с небольшим средним поперечным импульсом k-T: от 0,15 до 0,25 ГэВ на c. Эти условия указаны на самом рисунке.

Посмотрим на область малых относительных импульсов. Пионы каждой пары имеют одинаковый заряд и отталкиваются. Это меняет частоту регистрации близких импульсов одновременно с квантовым обменным усилением, которое мы уже разобрали.

Поэтому линии через точки получены с учётом кулоновского отталкивания — об этом говорит подпись внизу. Одной свободной гауссовой функции здесь недостаточно. Мы не измеряем радиус линейкой на рисунке: параметры извлекают, описывая всю форму корреляции физически обоснованной функцией.

\textbf{Переход:} Чтобы понять смысл этих параметров, уточним, всю ли систему характеризует пара с выбранным средним импульсом.

\textbf{Справка:} А. Краева (STAR), Universe 10, 188 (2024), рис. 4, §2.1, формула (2). Au+Au, √s\_NN=3 ГэВ, 0–10\%, 0.15<k\_T<0.25 ГэВ/c. Перед нами проекции одной трёхмерной функции: остальные компоненты q интегрированы по ±0.05 ГэВ/c, а не фиксированы в нуле. Положительные и отрицательные пары анализируются раздельно. Использована процедура Боулера–Синюкова:

\[C(\mathbf q)=\mathcal N\big[(1-\lambda)+\lambda K_{\rm C}(\mathbf q)(1+G(\mathbf q))\big].\]

K\_C — квадрат модуля кулоновской волновой функции, усреднённый по сферическому гауссову источнику (в статье радиус 5 фм); G — гауссова обменная часть с допустимыми смешанными членами. Сильное взаимодействие пионов этой формулой не моделируется. Для других пар сильное взаимодействие может существенно менять корреляцию. Цвета и обозначение STAR Preliminary оригинала сохранены.
